\def\topic{Farmer's calculus}
\ifdefined\MyHeaderIncluded
\else
\def\qwertyflag{true}

\ifdefined\MyHeaderIncluded

\else
  \newcommand{\foobar}{FooBar}
  \def\MyHeaderIncluded{true}
% if is till the end of the file

\documentclass[10pt,a4paper]{article}
\usepackage[utf8]{inputenc}
%\usepackage[german]{babel}
%\usepackage[cyrillic]{babel}
\usepackage[T1]{fontenc}
%\usepackage[T2A]{fontenc}
%\usepackage[english, russian]{babel}
\usepackage{amsmath}
\usepackage{amsfonts}
\usepackage{amssymb}
\usepackage{amsthm}
\usepackage{mathtools} %coloneqq

\usepackage{bussproofs} %

%\renewcommand{\thesubsection}{\arabic{subsection}}
%\newtheorem{thm}{Theorem}[section]
%\renewcommand{\thethm}{\arabic{subsection}.\arabic{thm}}

\usepackage{thmtools}
\usepackage{hyperref}
\hypersetup{
    colorlinks=true,
    linkcolor=blue,
    filecolor=magenta,      
    urlcolor=cyan,
    pdftitle={Overleaf Example},
    %pdfpagemode=FullScreen,
    }
%\usepackage{stix-type1}%\turnediota
\theoremstyle{definition}

\newtheorem{constructor}{Constructor}[section]
\newtheorem{iaxiom}{Axiom}
\newtheorem{axiom}{Axiom}
\newtheorem*{axiom*}{Axiom}
\newtheorem{definition}{Definition}[section]
\newtheorem*{definition*}{Definition}
\newtheorem{remark}{Remark}[section]
\newtheorem*{remark*}{Remark}
\newtheorem{recall}{Recall}[section]
\newtheorem*{recall*}{Recall}

\newtheorem{lemma}{Lemma}[section]
\newtheorem{theorem}{Theorem}[section]
\newtheorem*{theorem*}{Theorem}
\newtheorem{note}{Note}[section]
\usepackage{etoolbox}

%\newcommand*{\@Dashv}[2]{%
%  \reflectbox{$\m@th#1#2$}%
%}
\newcommand*{\separation}[3]{% 1)binder, 2)set, 3)condition
  \{#1\in #2: #3\}
}

% https://www.freecodecamp.org/news/simple-regex-tricks-for-beginners-3acb3fa257cb/
%  \\mathcal\{T\} |---> \topology
%  \\mathcal\{([^\}]*)\}   |--->  $1
\newcommand{\AND}{\ \ \&\ \ }
\newcommand{\any}{\mathrm{\ any}}
\newcommand{\borel}{\mathcal{B}}
\newcommand{\can}{\mathrm{can}}
\newcommand{\canonize}{\mathrm{canonize}}
\newcommand{\closure}{\mbox{cl}}
\newcommand{\complex}{\mathbb{C}}
\newcommand{\converges}{\longrightarrow}
\newcommand{\countable}{\mbox{\ countable}}
\newcommand{\ldefined}{\mbox{\ defined}}
\newcommand{\dom}{\mbox{dom}}
\newcommand{\Dom}{\mbox{Dom}}
\newcommand{\Done}{\noindent \textbf{Done.}\\}
\newcommand{\lexists}{\mbox{\ exists}}
\newcommand{\lforall}{\mbox{\ forall}}
\newcommand{\lunique}{\mbox{\ unique}}
\newcommand{\family}{\mathrm{\ family}}
\newcommand{\Fm}{\mathrm{Fm}}
\newcommand{\FuncSymb}{\mathrm{FuncSymb}}
\newcommand{\function}{\mathrm{\ function}}
\newcommand{\functional}{\mbox{\ functional}}
\newcommand{\IFF}{\ \Longleftrightarrow\ }
\renewcommand{\iff}{\ \longleftrightarrow\ }
\newcommand{\IMPLIES}{\ \Longrightarrow\ }
\renewcommand{\implies}{\ \longrightarrow\ }
\newcommand{\Inf}{\mathrm{Inf}}
\newcommand{\injective}{\mbox{\ injective}}
\newcommand{\injection}{\mbox{\ injection}}
\newcommand{\interior}{\mbox{int}}
\newcommand{\intersection}{\bigcap}
\newcommand{\isempty}{\mathrm{\ empty}}
\newcommand{\LB}{\mathrm{LB}}
\newcommand{\nghbset}{\mathcal{N}}
\newcommand{\nonempty}{\mathrm{\ nonempty}}
\newcommand{\Ord}{\mathrm{Ord}}
\newcommand{\power}{\mathcal{P}}
\newcommand{\pr}{\mathrm{pr}}
\newcommand{\fst}{\mathrm{fst}}
\newcommand{\snd}{\mathrm{snd}}
\newcommand{\PredSymb}{\mathrm{PredSymb}}
\newcommand{\Proof}{\noindent \textbf{Proof.}\\}
\newcommand{\ran}{\mathrm{ran}}
\newcommand{\Ran}{\mathrm{Ran}}
\newcommand{\real}{\mathbb{R}}
\newcommand{\relation}{\mathrm{\ relation}}
\newcommand{\set}{\mathrm{\ set}}
\newcommand{\singleton}{\mathrm{\ singleton}}
\newcommand{\semantics}{\mathrm{semantics}}
\newcommand{\semialgebra}{\mathcal{S}}
\newcommand{\Sup}{\mathrm{Sup}}
\newcommand{\surjective}{\mbox{\ surjective\ on\ }}
\newcommand{\Term}{\mathrm{Term}}
\newcommand{\Type}{\mathrm{Type}}
\newcommand{\topology}{\mathcal{T}}
\newcommand{\total}{\mbox{\ total\ on\ }}
\newcommand{\transitive}{\mathrm{\ transitive}}
\newcommand{\union}{\bigcup}
\newcommand{\lTrue}{\mathrm{True}}
\newcommand{\lFalse}{\mathrm{False}}
\newcommand{\lUncertain}{\mathrm{Uncertain}}
\newcommand{\lif}{\mathrm{if}}
\newcommand{\lthen}{\mathrm{then}}
\newcommand{\lelse}{\mathrm{else}}
\newcommand{\lexception}{\mathrm{exception}}
\newcommand{\lundef}{\mathrm{undef}}
\newcommand{\univocal}{\mathrm{\ univocal}}
\newcommand{\univocalat}{\mathrm{\ univocal\ at\ }}
\newcommand{\UB}{\mathrm{UB}}
\newcommand{\Un}{\mathcal{U}}
\newcommand{\Var}{\mathrm{Var}}

\newcommand{\notvarphi}{\phi}
\renewcommand{\phi}{\varphi}

\newcommand{\comment}[1]{\quad\mbox{(#1)}}
% local commands
%\preto\section{\undef\blipo} %etoolbox
\preto\subsection{\undef\blipo} %etoolbox
% https://tex.stackexchange.com/questions/105684/limit-macro-newcommand-scope-to-section
\usepackage{graphicx}
\usepackage{authblk}

\setlength\parindent{0pt} % whole doc of \noindent
\usepackage{soul}
%\usepackage{mathgl}

\fi % guard
\ifdefined\topic
\else
\def\topic{No Titel}
\fi
\title{\topic}
\author[1]{Georgy Dunaev}
\begin{document}
\maketitle
\fi

%\section{Appendix 1: Alternative partial function semantics [unrelated to the rest of the article]}
\section{Partial function semantics following W.\,Farmer}

This section uses partial first-order logic (PFOL). The notation
$t\downarrow$ asserts that $t$ denotes, and $s\simeq t$ denotes Kleene
equality: the terms either denote the same value or are both undefined.
Variables denote; compound terms may fail to denote. Function evaluation
is defined by a unique description, independently of the total union-based
evaluation used in other sections.

\begin{note}[Three conventions for an absent value]
The newer Obelisk catalog's \texttt{DefValueAtArg} (3.25) uses a total
conditional evaluation: it returns $\V$ when there is no represented
output. Plain union evaluation, used in another local construction of
this book, returns $\union\emptyclass=\emptyclass$ in that case.
PFOL instead leaves the unique description undefined. These are three
distinct conventions. For the empty relation, for example, the value at
any argument is respectively $\V$, $\emptyclass$, or undefined.
The definedness assertions in this section belong to the last convention.
\end{note}

\subsection{A proof in PFOL}
The argument follows the notation of William M. Farmer and Joshua D.
Guttman, ``A Set Theory with Support for Partial Functions''. The PFOL
axiom numbers below refer to that presentation, especially Section~2.3;
its guarded application constructor appears in Section~3.2.
\footnote{William M. Farmer and Joshua D. Guttman,
``A Set Theory with Support for Partial Functions'',
\emph{Studia Logica} 66 (2000), 59--78.
\url{https://web.cs.wpi.edu/~guttman/pubs/set_theory_partial_fns.pdf}.}
For the set-theoretic statements below, indices and output variables are
sets, ordered pairs use the usual injective encoding, and relations may
be classes. This convention avoids treating a proper class as a component
of a set-coded pair.
\begin{definition}[value of a function]
\[
f(x) \simeq \left(\iota y, \langle x,y\rangle\in f\right)
\]
\end{definition}
Here the descriptor tests uniqueness at the particular input $x$.
Farmer--Guttman's application constructor instead includes the additional
conjunct that $f$ is a function. The two agree for functions; they can
differ for other relations. For example,
\[
r=\{\langle0,0\rangle,\langle1,0\rangle,\langle1,1\rangle\}
\]
has local value $r(0)=0$ and no local value at $1$, while the globally
guarded constructor is undefined at every input because $r$ is not
univocal. The local convention is intentional in the following results.

\begin{theorem}[Evaluation of a univocal relation]
If $f$ is univocal (at most one output for each input), and $x,y$ are
variables, then
\[
\left(f(x)\simeq y\right)\iff \left(\langle x,y\rangle\in f\right)
\]
\end{theorem}
\Proof
($\Leftarrow$) We use the following form PFOL$10'$ of PFOL10:
\[
(\iota x,\ \varphi)\downarrow
\quad\overset{\mbox{PFOL10}'}{\Longleftrightarrow}\quad
(\exists!x,\ \varphi)
\]
So the proof is the following:
\[
\left(\langle x,y\rangle\in f\right)
\quad\overset{\exists\mbox{--intro}}{\Longrightarrow}\quad
\left(\exists w,\ \langle x,w\rangle\in f\right)
\]
\[
f\mbox{ univocal}
\quad\overset{\mbox{def}}{\Longrightarrow}\quad
\left(! w,\ \langle x,w\rangle\in f\right)
\]
\[
\left(\exists w,\ \langle x,w\rangle\in f\right)
\LAND
\left(! w,\ \langle x,w\rangle\in f\right)
\quad\overset{\exists!\text{--intro}}{\Longrightarrow}\quad
\left(\exists! w,\ \langle x,w\rangle\in f\right)
\]
\[
\left(\exists! w,\ \langle x,w\rangle\in f\right)
\quad\overset{\mbox{PFOL10}'}{\Longrightarrow}\quad
\left(\iota w,\ \langle x,w\rangle\in f\right)\downarrow
\]
\[
\left(\iota w,\ \langle x,w\rangle\in f\right)\downarrow
\quad\overset{\mbox{PFOL11}}{\Longrightarrow}\quad
\langle x,(\iota w,\ \langle x,w\rangle\in f)\rangle\in f
\]
\[
\langle x,(\iota w,\ \langle x,w\rangle\in f)\rangle\in f
\LAND
\left(\langle x,y\rangle\in f\right)
\quad\overset{f \mbox{ univocal}}{\Longrightarrow}\quad
(\iota w,\ \langle x,w\rangle\in f) \simeq y
\]
\[
f(x) \simeq y
\]
%\left(! w,\ \langle x,w\rangle\in f\right)
%\[
%\left(\iota w,\ \langle x,w\rangle\in f\right)\downarrow

%\]
\noindent($\Rightarrow$)
\[(y\downarrow)\quad\mbox{(by PFOL8)}\]
\[
(y\downarrow)\LAND((\iota w,\ \langle x,w\rangle\in f)\simeq y) \quad\overset{\mbox{PFOL7}}{\Longrightarrow}\quad
(\iota w,\ \langle x,w\rangle\in f)\downarrow \]
\[
\left(\iota w,\ \langle x,w\rangle\in f\right)\downarrow
\quad\overset{\mbox{PFOL11}}{\Longrightarrow}\quad
\langle x,(\iota w,\ \langle x,w\rangle\in f)\rangle\in f
\]
\[
\langle x,(\iota w,\ \langle x,w\rangle\in f)\rangle\in f
\LAND
\left((\iota w,\ \langle x,w\rangle\in f)\simeq y\right)
\quad\overset{\mbox{PFOL7}}{\Longrightarrow}\quad
\left(\langle x, y\rangle\in f\right)
\]
\\
\Done
\begin{theorem}[Domain of unique-description evaluation]
\[
f(x)\downarrow \iff x\in\cc{x}{\exists!y,\ \langle x,y\rangle\in f}
\]
\end{theorem}
\Proof
By the definition of $f(x)$ and PFOL$10'$, it denotes exactly when
$\exists!y,\ \langle x,y\rangle\in f$. The displayed class collects
precisely these inputs, giving the claimed equivalence.
\Done
\begin{theorem}[The graph retained by local evaluation]
For any class relation $r$ of ordered pairs of sets, define
\[
U(r)=\{\langle x,y\rangle\in r:
       \forall z\in V\ (\langle x,z\rangle\in r\to z=y)\}.
\]
Then $U(r)$ is univocal, $U(U(r))=U(r)$, and
$U(r)(x)\simeq r(x)$ for every set $x$. Two relations have the same
local evaluations at every set input if and only if their $U$-graphs
are equal.
\end{theorem}
\Proof
At an input with exactly one output, $U$ retains that output. At an
input with no output or with at least two outputs, it retains none.
This proves univocality and equality of the partial evaluations. Applying
$U$ again changes no singleton or empty fibre, proving idempotence.
If the retained graphs are equal, equality of local evaluations follows.
Conversely, if all local evaluations agree, each retained pair on one
side supplies a denoting value on both sides and hence is retained on
the other. Extensionality gives equality of the retained graphs.
\Done

This operation is not monotone for inclusion of relations. With
$0\ne1$, let $r=\{\langle0,0\rangle\}$ and
$s=r\cup\{\langle0,1\rangle\}$. Then $r\subseteq s$, but
$U(r)=r$ and $U(s)=\emptyclass$. Adding a possible output can destroy
local definedness. Consequently the extensionality theorem below
requires univocal graphs: evaluation intentionally forgets every
ambiguous fibre of a general relation.

\subsection{Equality and composition of partial evaluations}
Strict equality $s=t$ is true only when both terms denote the same value.
Kleene equality therefore satisfies
\[
s\simeq t\iff (s=t)\lor(s\uparrow\land t\uparrow),
\qquad t=t\iff t\downarrow,
\]
where $t\uparrow$ abbreviates $\neg(t\downarrow)$. These identities follow
by separating the cases in which both, one, or neither term denotes.
Consequently the denoting-variable hypothesis in the evaluation theorem
matters: if $f(x)$ and a term $t$ are both undefined, $f(x)\simeq t$ is
true, while strict pair membership $\langle x,t\rangle\in f$ is false.

\begin{theorem}[Extensionality of partial function graphs]
For univocal class relations $f,g$ of ordered pairs of sets,
\[
f=g\iff \forall x\in V,\quad f(x)\simeq g(x).
\]
\end{theorem}
\Proof
Equal relations have the same fibres, so their evaluations either both
fail to denote or denote the same unique member. Conversely, suppose
all their evaluations are Kleene equal. If $\langle x,y\rangle\in f$,
then $f(x)$ denotes $y$. Kleene equality forces $g(x)$ to denote the same
$y$, so the evaluation theorem gives $\langle x,y\rangle\in g$.
This proves $f\subseteq g$; exchanging the relations proves the reverse
inclusion. Class extensionality gives equality.
\Done

\begin{theorem}[Composition and its domain]
Let $f,g$ be univocal class relations of ordered pairs of sets, and put
\[
h=g\circ f=\{\langle x,z\rangle:
       \exists y\in V\ (\langle x,y\rangle\in f
                         \land\langle y,z\rangle\in g)\}.
\]
Then $h$ is univocal, and for every set $x$,
\[
h(x)\simeq g(f(x)),
\qquad
h(x)\downarrow\iff
\exists y\in V\ (\langle x,y\rangle\in f\land y\in\dom(g)).
\]
\end{theorem}
\Proof
Two outputs of $h$ at $x$ have intermediate witnesses $y,y'$. Univocality
of $f$ gives $y=y'$, and then univocality of $g$ identifies the outputs.
If $f$ has no output at $x$, the composite graph has none; also $f(x)$
is undefined. Strict pair membership makes every witness condition in
$g(f(x))$ false, so this term is undefined too. If $f(x)$ denotes $y$,
the fibre of $h$ at $x$ is exactly the fibre of $g$ at $y$. It is either
empty on both sides or has the same unique output on both sides. This
proves the Kleene equality and the stated domain criterion.
\Done

Univocality cannot be omitted from the composition theorem. Take
\[
f=\{\langle0,0\rangle,\langle0,1\rangle\},\qquad
 g=\{\langle0,0\rangle,\langle1,0\rangle\}.
\]
Then $g\circ f=\{\langle0,0\rangle\}$, whose value at $0$ is $0$.
But $f(0)$ is undefined under unique-description evaluation, so
$g(f(0))$ is undefined. A multivalued intermediate relation can lose
its ambiguity after composition; local uniqueness does not commute with
that operation.

\subsection{Extension by definitions}
The following formulas describe how a descriptor can be introduced under
an existence hypothesis. They are derived schemata, not further PFOL axioms.
In particular, unique existence must be asserted under the domain
hypothesis: placing that hypothesis inside an implication governed by
$\exists!y$ would impose uniqueness even when the hypothesis is false.
\pagebreak
The intended conditional unique-existence premise is
\[
\forall f \forall x,\ ((x\in \dom(f))\land (f\mbox{ univocal at }x)) \to (\exists! y,\ (x,y)\in f)
\]
from which the intended evaluation rule is
\[
\forall f \forall x,\ ((x\in \dom(f))\land (f\mbox{ univocal at }x)) \to (x,f(x))\in f
\]

A total extension by a function symbol requires a unique value for every
input. A partial extension instead retains the domain hypothesis above.
The distinction is the reason for using partial descriptions in this section.

In general, a closed formula
\[
\forall \vec{x} ,\ \exists! y,\ \varphi
\]
implies
\[
\forall \vec{x} ,\ \varphi[(\iota y,\ \varphi)/y]
\]
.
\iffalse
IMHO, the most important part here is that y is not free in the context below.
(the reason is identical to the one in the Generalization Rule)
YES!
Important: get rid if quantifier over arguments.
If the context doesn't have y, the substitution of "whole phi" would be already done.

something that a lil bit concerning is that context is not seen inside I, but it may not be important.


So
\fi

The first rule below is admissible by PFOL10 and PFOL11: unique
existence yields definedness of the descriptor, and definedness yields
its witness formula. The second additionally requires a choice descriptor
with its own witness axiom. In both rules, substitution must be free for
$y$, with bound variables renamed where necessary. No generalization is
performed, so $y\notin\mathrm{FV}(\Gamma)$ is not needed for these
witness rules. Freshness conditions for a newly introduced symbol or an
eigenvariable belong to different rules.
\begin{prooftree}
\AxiomC{$\Gamma \vdash \exists! y,\ \varphi$}
%\RightLabel{Assume/Ident}
%\UnaryInfC{$
%P_3,P_2,P_1 \vdash R
%$}
\RightLabel{Unique witness}
\UnaryInfC{$
\Gamma \vdash \varphi[(\iota y,\ \varphi)/y]
$}
%\RightLabel{by }
\end{prooftree}
and
\begin{prooftree}
\AxiomC{$\Gamma \vdash \exists y,\ \varphi$}
\RightLabel{Choice witness}
\UnaryInfC{$
\Gamma \vdash \varphi[(\epsilon y,\ \varphi)/y]
$}
\end{prooftree}
\begin{theorem}[A guarded partial definition]
Let $A(\vec x)$ and $\varphi(\vec x,y)$ be old-language formulas whose
free variables are among those displayed, with $y$ not free in $A$.
Suppose a theory $T$ entails
\[
\forall\vec x\ (A(\vec x)\to\exists!y\ \varphi(\vec x,y)).
\]
Assume $\vec x$ is a nonempty tuple of distinct variables. Add a new
partial operator $k$ of that arity with the axioms
\[
\forall\vec x\ (k(\vec x)\downarrow\iff A(\vec x)),\qquad
\forall\vec x\ (A(\vec x)\to\varphi(\vec x,k(\vec x))).
\]
Every model of $T$ has exactly one expansion interpreting $k$ subject
to these axioms. Its value is described by
\[
k(\vec x)\simeq
\bigl(\iota y,\ A(\vec x)\land\varphi(\vec x,y)\bigr).
\]
The extension is conservative for old-language semantic consequences.
\end{theorem}
\Proof
Fix a model of $T$ and a tuple of denoting arguments. If $A$ is true,
there is exactly one witness for $\varphi$; assign that witness as the
value of $k$. If $A$ is false, leave $k$ undefined. This defines a
partial operation and satisfies both axioms. Any operation satisfying
them must be undefined in the second case and equal the unique witness
in the first, proving uniqueness. The displayed descriptor has exactly
this behavior: its predicate has one witness when $A$ holds and no
witness when $A$ fails.

An old-language sentence has the same truth value in a model and its
expansion. Thus every old-language countermodel to a proposed consequence
of $T$ expands to a countermodel of the extension. No new old-language
semantic consequence is obtained. This is a model argument; a claim
about derivability additionally uses completeness of the chosen calculus.
\Done

The domain axiom is essential. A guarded witness axiom alone allows
arbitrary values away from $A$ and does not determine a unique partial
operation. Likewise, inserting $A\to\varphi$ instead of
$A\land\varphi$ inside the descriptor gives the wrong off-domain
predicate: when $A$ is false, every object is a witness. In a one-object
model the resulting descriptor even denotes there. Finally, if one
substitutes possibly undefined terms for $\vec x$, strict operator
application remains undefined when an argument is undefined; the theorem
and its displayed defining equations quantify over denoting variables.


\pagebreak




%\ifdefined\partialfunctionssemanticsflag
%\undef\partialfunctionssemanticsflag
%\end{document}
%\fi

\ifdefined\qwertyflag
\undef\qwertyflag
\end{document}
\fi
